\documentclass[11pt,a4paper]{article}

\usepackage[a4paper,margin=2.5cm]{geometry}
\usepackage{fontspec}
\usepackage{titlesec}
\usepackage{longtable}
\usepackage{booktabs}
\usepackage{array}
\usepackage{ragged2e}
\usepackage{xcolor}
\usepackage{hyperref}
\usepackage{enumitem}
\usepackage{fancyhdr}

\setmainfont{Georgia}
\definecolor{brandgreen}{HTML}{1F6F5C}
\definecolor{rowgray}{HTML}{F8F9FA}

\hypersetup{
  colorlinks=true,
  linkcolor=brandgreen,
  urlcolor=brandgreen,
  pdftitle={Anforderungskatalog -- LLM-Driven Climate-Health ETL Platform},
  pdfauthor={Sergi Koniashvili}
}

\titleformat{\section}{\large\bfseries\color{brandgreen}}{\thesection}{0.5em}{}
\titleformat{\subsection}{\normalsize\bfseries}{\thesubsection}{0.5em}{}

\setlength{\headheight}{14pt}
\pagestyle{fancy}
\fancyhf{}
\fancyhead[L]{LLM-Driven Climate-Health ETL Platform}
\fancyhead[R]{Anforderungskatalog / Requirements Catalog}
\fancyfoot[C]{\thepage}

\setlength{\emergencystretch}{3em}
\setlength{\tabcolsep}{4pt}
\renewcommand{\arraystretch}{1.3}
\newcolumntype{L}[1]{>{\RaggedRight\arraybackslash}p{#1}}

\begin{document}

\begin{titlepage}
\centering
\vspace*{2.5cm}
{\Huge\bfseries\color{brandgreen} Anforderungskatalog}\\[0.3cm]
{\Large Requirements Catalog}\\[1.5cm]
{\LARGE An LLM-Driven ETL Platform for\\ Climate-Health Data Integration}\\[1.8cm]
{\large Bachelor-Arbeit / Bachelor Thesis}\\[0.3cm]
{\large Universit\"at Bremen -- Fachbereich Informatik}\\[1.8cm]
{\large Autor / Author: \textbf{Sergi Koniashvili}}\\[0.4cm]
{\large Stand / Date: \today}\\[2.5cm]
{\normalsize Dieses Dokument beschreibt den umfassenden funktionalen und nicht-funktionalen\\
Anforderungskatalog f\"ur die entwickelte Software-Plattform, einschlie\ss{}lich aller\\
Erweiterungen (Copilot, Literatur-RAG, Demografie, statistische Lags und Downscaling)\\
sowie den aktuellen Umsetzungsstand.}
\vfill
\end{titlepage}

\tableofcontents
\newpage

%=====================================================================
\part*{Teil I -- Deutsch}
\addcontentsline{toc}{part}{Teil I -- Deutsch}

\section{Einleitung und Zielsetzung}

Viele Infektionskrankheiten sind hochgradig klimasensitiv: Ihre biologischen Vektoren
(z.\,B. \emph{Aedes aegypti}- und \emph{Anopheles}-M\"ucken) reagieren empfindlich
auf Temperatur, Niederschlag und Luftfeuchtigkeit. Hochaufgel\"oste Klimadaten dienen
daher als essenzieller Fr\"uhindikator f\"ur Ausbruchsrisiken. 

Die im Rahmen dieser Bachelorarbeit entwickelte Plattform erm\"oglicht es
Forschenden, Beh\"orden und Interessierten ohne Programmier- oder Klimawissenschafts\-kenntnisse,
per nat\"urlichsprachlicher Eingabe oder interaktiver Steuerung analysebereite,
klimabezogene Gesundheits- und Demografiedatens\"atze zu erstellen. Ein Large Language
Model (LLM) \"ubersetzt Anfragen in strukturierte Integrationspl\"ane, ruft Daten aus
offiziellen wissenschaftlichen Archiven ab, harmonisiert unterschiedliche zeitliche
und r\"aumliche Aufl\"osungen und liefert eine fundierte, laiengerechte Erkl\"arung
sowie statistische Korrelationsanalysen.

\textbf{Validierungsszenario:} Denguefieber in Thailand, mit voller Generalisierbarkeit
auf weitere vektor\"ubertragene Infektionskrankheiten (Malaria, Cholera, West-Nil-Virus
u.\,a.) und globale Regionen.

\section{Funktionale Anforderungen (FR)}

\begin{longtable}{@{}>{\bfseries}p{1.6cm} L{7.4cm} L{3.6cm} L{2.4cm}@{}}
\toprule
ID & Beschreibung & Herkunft & Status \\
\midrule
\endhead
FR-01 & Das System muss eine nat\"urlichsprachliche Anfrage (Krankheit, Region, Klimavariablen, Zeitspanne) entgegennehmen und per LLM in einen strukturierten Integrationsplan \"ubersetzen. & Projektkonzept & Umgesetzt \\
FR-02 & Die Welt muss in ausw\"ahlbare Regionen unterteilt sein; die Nutzerin/der Nutzer muss Ergebnisse regionsspezifisch abrufen k\"onnen (39 L\"ander hinterlegt). & Auflage Betreuer & Umgesetzt \\
FR-03 & Das System muss jeden Verarbeitungsschritt (Abruf, Validierung, Harmonisierung, Verkn\"upfung) sichtbar protokollieren. & Auflage Betreuer & Umgesetzt \\
FR-04 & Klimadaten m\"ussen ausschlie\ss{}lich aus wissenschaftlichen/offiziellen Quellen stammen (ERA5, NASA POWER, TMD). & Auflage Betreuer & Umgesetzt \\
FR-05 & Die Nutzerin/der Nutzer muss zwischen mehreren wissenschaftlichen Datenquellen w\"ahlen k\"onnen. & Auflage Betreuer & Umgesetzt \\
FR-06 & Die Nutzerin/der Nutzer muss eine eigene Datenquelle per URL angeben k\"onnen (Fallzahlen, Wetter, Demografie). & Auflage Betreuer & Umgesetzt \\
FR-07 & Die Nutzerin/der Nutzer muss eine eigene Datenquelle als CSV-Datei hochladen k\"onnen (Fallzahlen, Wetter, Demografie). & Auflage Betreuer & Umgesetzt \\
FR-08 & Empfangene und erzeugte Datens\"atze m\"ussen als CSV und JSON heruntergeladen werden k\"onnen. & Auflage Betreuer & Umgesetzt \\
FR-09 & Das System muss jeden Schritt und die epidemiologische Beziehung in verst\"andlicher Sprache erkl\"aren. & Projektkonzept & Umgesetzt (LLM-generiert) \\
FR-10 & Das System muss Datenabruf und Validierung unter Beachtung von API-Ratenlimits durchf\"uhren, ohne die UI zu blockieren. & Nutzer\-anforderung & Umgesetzt \\
FR-11 & Das System muss bereits verarbeitete Daten cachen und mit Zeitstempel der letzten \"Uberpr\"ufung (\texttt{last\_verified}) anzeigen. & Nutzer\-anforderung & Umgesetzt \\
FR-12 & Das System muss modular auf weitere Krankheiten und Regionen erweiterbar sein, ohne die Kernarchitektur zu \"andern. & Projektkonzept & Umgesetzt \\
FR-13 & Das System muss einen interaktiven Copilot-Chatassistenten bieten, der Fachfragen beantwortet, Formularparameter vorschl\"agt und Suchl\"aufe direkt ausl\"ost. & Erweiterung / UX & Umgesetzt \\
FR-14 & Das System muss eine wissenschaftliche Literatur-Wissensbasis (WHO, Lancet, IPCC, Mordecai et al.) bereitstellen, Vergleiche erl\"autern und RAG-In-Context-Grounding erm\"oglichen. & Fachliche Anforderung & Umgesetzt \\
FR-15 & Das System muss Metadaten aus wissenschaftlichen PDF-Dateien oder Web-Links automatisiert per KI extrahieren und strukturieren. & Erweiterung / KI & Umgesetzt \\
FR-16 & Das System muss eine Live-Suche f\"ur externe Fallzahldatenquellen (WHO GHO Indikatoren und HDX-Datens\"atze) bereitstellen. & Erweiterung / Daten & Umgesetzt \\
FR-17 & Das System muss offizielle Demografiedaten (World Bank Indikatoren, UN WPP, Upload) durchsuchbar anbinden und Inzidenzraten pro 100.000 Einwohner berechnen. & Fachliche Anforderung & Umgesetzt \\
FR-18 & Das System muss f\"ur benutzerdefinierte Uploads einen l\"uckenlosen Transformations-Audit-Trail f\"uhren und verst\"andlich erl\"autern. & Qualit\"ats\-anforderung & Umgesetzt \\
FR-19 & Das System muss zeitverz\"ogerte Kreuzkorrelationen (Lags 0--3) mit Pearson-$r$ und $p$-Wert-Signifikanz berechnen. & Analytische Anforderung & Umgesetzt \\
FR-20 & Das System muss bei t\"aglicher Aufl\"osung eine massenerhaltende glatte zeitliche Interpolation durchf\"uhren und Aufl\"osungshinweise geben. & Methodische Anforderung & Umgesetzt \\
FR-21 & Das System muss eine vollst\"andig dreisprachige Benutzeroberfl\"ache (Deutsch, Englisch, Georgisch) mit Live-Umschaltung bieten. & Nutzer\-anforderung & Umgesetzt \\
FR-22 & Das System muss einen druckreifen wissenschaftlichen PDF-Analysebericht mit Diagrammen, Tabellen und Zitationen generieren k\"onnen. & Nutzer\-anforderung & Umgesetzt \\
FR-23 & Das System muss eine interaktive Weltkarte zur intuitiven visuellen L\"anderauswahl bereitstellen. & Usability-Anforderung & Umgesetzt \\
\bottomrule
\end{longtable}

\section{Nicht-funktionale Anforderungen (NFR)}

\begin{longtable}{@{}>{\bfseries}p{1.6cm} L{7.4cm} L{3.6cm} L{2.4cm}@{}}
\toprule
ID & Beschreibung & Herkunft & Status \\
\midrule
\endhead
NFR-01 & Ratenlimit-Schutz: Das System darf das Kontingent der LLM-API nicht \"uberschreiten (Sliding-Window-Begrenzung). & Betriebs\-anforderung & Umgesetzt \\
NFR-02 & Wissenschaftliche Zitierf\"ahigkeit: Jede Antwort muss verwendete Prim\"arquellen (Klima, F\"alle, Demografie, Literatur) exakt zitieren. & Wissenschafts\-praxis & Umgesetzt \\
NFR-03 & Portabilit\"at: Die gesamte Plattform (FastAPI-Backend + Flutter-Web-Frontend) muss als einzelner Docker-Container lauff\"ahig sein. & DevOps / Betrieb & Umgesetzt \\
NFR-04 & Datensicherheit: API-Schl\"ussel und Zugangsdaten d\"urfen niemals im Quellcode oder im Client-Bundle exponiert werden. & Sicherheits\-anforderung & Umgesetzt (\texttt{.env}, \texttt{.gitignore}) \\
NFR-05 & Verst\"andlichkeit: Erkl\"arungen m\"ussen f\"ur Personen ohne Vorkenntnisse in Klimamodellierung oder Informatik intuitiv begreifbar sein. & Projektkonzept & Umgesetzt \\
NFR-06 & Wartbarkeit: Strikte Trennung von Datenextraktion, Harmonisierung, statistischer Analyse und Pr\"asentation in modularen Services. & Software-Engineering & Umgesetzt \\
NFR-07 & RAG-Integrit\"at: Das System nutzt In-Context Grounding statt Modell-Feintuning; neuronale Gewichte bleiben unver\"andert. & Wissenschafts\-ethik & Umgesetzt \\
NFR-08 & Mathematische Massenerhaltung: Die zeitliche Interpolation auf Tagesebene erh\"alt exakt die Gesamtsumme der gemeldeten Prim\"arf\"alle. & Datenintegrit\"at & Umgesetzt \\
NFR-09 & Sprachkonsistenz: Vermeidung von Sprachmischungen; alle Interface-Texte liegen konsistent in de, en und ka vor. & Usability / i18n & Umgesetzt \\
\bottomrule
\end{longtable}

\section{Architektur\"uberblick}

Die Plattform basiert auf einer Drei-Schichten-Architektur:
\begin{enumerate}[label=(\arabic*)]\sloppy
  \item \textbf{Frontend (Flutter Web):} Reaktive Einzelseiten-Applikation mit mehrsprachiger Benutzeroberfl\"ache, Copilot-Drawer, Literatur-Korpus-Verwaltung, Daten-Audit-Anzeige, interaktiver Weltkarte und PDF-Export.
  \item \textbf{Backend (Python, FastAPI):} Modulare Service-Architektur f\"ur Pipeline-Orchestrierung (\texttt{etl}), Klimadatenzugriff (\texttt{climate}, \texttt{tmd}), Surveillance-Erfassung (\texttt{case\_data}, \texttt{who\_gho}, \texttt{hdx}), Demografie (\texttt{population}), Korrelationsstatistik (\texttt{statistics}), automatisierte Literatur-Extraktion (\texttt{literature\_extractor}) sowie LLM-Copilot-Dienste (\texttt{llm}, \texttt{chat}).
  \item \textbf{Wissenschaftliche Prim\"arquellen:} Reanalyse-APIs (Open-Meteo ERA5, NASA POWER, TMD), globale Gesundheitsdaten (OpenDengue, WHO GHO, HDX), Demografie (Weltbank WDI, UN WPP) und Google Gemini LLM.
\end{enumerate}

\section{Validierung}

Das Pipeline-Design wurde prim\"ar anhand von Denguefieber in Thailand (1990--2023)
validiert. Es ist vollst\"andig generalisierbar auf alle 39 hinterlegten L\"ander
sowie frei konfigurierbar f\"ur beliebige nutzerdefinierte Datens\"atze.

\newpage
%=====================================================================
\part*{Part II -- English}
\addcontentsline{toc}{part}{Part II -- English}

\section{Introduction and Objectives}

Many infectious diseases are strongly climate-sensitive: their biological vectors
(e.g., \emph{Aedes aegypti} and \emph{Anopheles} mosquitoes) exhibit acute sensitivity
to ambient temperature, rainfall, and relative humidity. High-resolution meteorological
data therefore serves as a vital early indicator of disease transmission suitability
and outbreak hazards.

The software platform developed for this bachelor thesis enables researchers, public
health officials, and decision-makers without programming or meteorological backgrounds
to obtain analysis-ready, climate-harmonized epidemiological datasets. Through
natural language input or direct form selection, a Large Language Model (LLM) translates
user intent into an executable ETL pipeline plan, queries verified international repositories,
harmonizes conflicting temporal and spatial resolutions, computes demographic incidence rates,
and performs statistical lag cross-correlations with plain-language explanations.

\textbf{Validation Scenario:} Dengue fever in Thailand, with architectural extensibility
to further vector-borne diseases (malaria, cholera, West Nile virus) and global regions.

\section{Functional Requirements (FR)}

\begin{longtable}{@{}>{\bfseries}p{1.6cm} L{7.4cm} L{3.6cm} L{2.4cm}@{}}
\toprule
ID & Description & Origin & Status \\
\midrule
\endhead
FR-01 & The system must accept natural language queries (disease, region, climate variables, dates) and translate them via LLM into a structured integration plan. & Project concept & Implemented \\
FR-02 & The world must be partitioned into selectable regions; users must be able to retrieve region-specific datasets (39 countries registered). & Supervisor requirement & Implemented \\
FR-03 & The system must visibly log every pipeline execution step (retrieval, validation, harmonization, joining, statistical evaluation). & Supervisor requirement & Implemented \\
FR-04 & Meteorological series must originate strictly from authoritative scientific sources (ERA5, NASA POWER, TMD). & Supervisor requirement & Implemented \\
FR-05 & Users must be able to choose between multiple verified scientific data sources. & Supervisor requirement & Implemented \\
FR-06 & Users must be able to supply custom data sources via web URL (cases, weather, demographics). & Supervisor requirement & Implemented \\
FR-07 & Users must be able to upload custom datasets as CSV files (cases, weather, demographics). & Supervisor requirement & Implemented \\
FR-08 & Harmonized datasets must be downloadable locally in standard CSV and JSON formats. & Supervisor requirement & Implemented \\
FR-09 & The system must explain pipeline transformations and climate-health associations in clear, plain language. & Project concept & Implemented (LLM-generated) \\
FR-10 & The system must respect upstream API rate limits asynchronously without blocking user interactions. & User requirement & Implemented \\
FR-11 & The system must cache processed datasets and surface a verification freshness timestamp (\texttt{last\_verified}). & User requirement & Implemented \\
FR-12 & The system must be modularly extensible to additional diseases and regions without architectural refactoring. & Project concept & Implemented \\
FR-13 & The system must provide an interactive Copilot chat assistant capable of domain dialogue, form prefill, and one-click execution. & UX requirement & Implemented \\
FR-14 & The system must incorporate a verified literature knowledge base (WHO, Lancet, IPCC, Mordecai et al.) with RAG in-context grounding and comparison. & Domain requirement & Implemented \\
FR-15 & The system must automatically extract and structure bibliographic metadata from uploaded PDFs or article URLs using AI. & AI requirement & Implemented \\
FR-16 & The system must support real-time online discovery of official surveillance sources (live query of WHO GHO indicators and HDX datasets). & Data requirement & Implemented \\
FR-17 & The system must integrate demographic indicators (World Bank Open Data, UN WPP, custom upload) and compute disease incidence per 100,000 population. & Domain requirement & Implemented \\
FR-18 & The system must maintain an end-to-end Data Transformation Audit trail for custom uploads with natural-language explanation. & Audit requirement & Implemented \\
FR-19 & The system must calculate time-lagged cross-correlations (Lags 0--3) with Pearson-$r$ and $p$-value statistical significance. & Analytical requirement & Implemented \\
FR-20 & The system must apply mass-preserving smooth temporal downscaling for daily surveillance resolution with user guidance. & Methodological requirement & Implemented \\
FR-21 & The system must offer a fully trilingual user interface (English, German, Georgian) with instant runtime language switching. & User requirement & Implemented \\
FR-22 & The system must generate downloadable, publication-ready scientific PDF report digests with charts, tables, and citations. & User requirement & Implemented \\
FR-23 & The system must feature an interactive world map for visual region and coordinate selection. & Usability requirement & Implemented \\
\bottomrule
\end{longtable}

\section{Non-Functional Requirements (NFR)}

\begin{longtable}{@{}>{\bfseries}p{1.6cm} L{7.4cm} L{3.6cm} L{2.4cm}@{}}
\toprule
ID & Description & Origin & Status \\
\midrule
\endhead
NFR-01 & Rate Limit Protection: The system must enforce sliding-window quotas to prevent exceeding upstream LLM API limits (15 req/min). & Operational requirement & Implemented \\
NFR-02 & Scientific Traceability: Every system response must provide rigorous academic citations for all primary sources used. & Academic standards & Implemented \\
NFR-03 & Portability: The complete software suite (FastAPI backend + Flutter Web) must be deployable as a single unified Docker container. & DevOps / Portability & Implemented \\
NFR-04 & Security: Credentials and private API keys must never be committed to source control or exposed to client-side bundles. & Security standards & Implemented (\texttt{.env}, \texttt{.gitignore}) \\
NFR-05 & Comprehensibility: Narratives and explanations must be interpretable by stakeholders without formal training in epidemiology or meteorology. & User-centric design & Implemented \\
NFR-06 & Maintainability: Clean separation of configuration, data harvesting, orchestration, analysis, and presentation into decoupled modules. & Software engineering & Implemented \\
NFR-07 & Grounding Integrity: RAG in-context grounding must be strictly maintained; model weights remain frozen, preventing hallucinated citations. & AI ethics & Implemented \\
NFR-08 & Mathematical Conservation: Daily downscaling must strictly conserve the integral case mass of the underlying official surveillance reports. & Data integrity & Implemented \\
NFR-09 & Linguistic Consistency: Full parity across German, English, and Georgian without mixed-language artifacts. & i18n standards & Implemented \\
\bottomrule
\end{longtable}

\section{Architecture Overview}

The platform follows a decoupled three-tier architecture:
\begin{enumerate}[label=(\arabic*)]\sloppy
  \item \textbf{Presentation Tier (Flutter Web):} Reactive single-page web client supporting trilingual localization, Copilot drawer, literature corpus browser, transformation audit view, interactive world map, and PDF export.
  \item \textbf{Service \& Orchestration Tier (Python, FastAPI):} Modulated services for pipeline execution (\texttt{etl}), meteorological querying (\texttt{climate}, \texttt{tmd}), disease surveillance (\texttt{case\_data}, \texttt{who\_gho}, \texttt{hdx}), demographics (\texttt{population}), lag statistics (\texttt{statistics}), literature parsing (\texttt{literature\_extractor}), and LLM copilot services (\texttt{llm}, \texttt{chat}).
  \item \textbf{Authoritative Upstream Tier:} Climate reanalysis archives (Open-Meteo ERA5, NASA POWER, TMD), surveillance repositories (OpenDengue, WHO GHO, HDX), demographic databases (World Bank WDI, UN WPP), and Google Gemini LLM.
\end{enumerate}

\section{Validation}

Pipeline efficacy is empirically demonstrated using historical dengue surveillance in
Thailand (1990--2023). Architectural generality ensures immediate applicability
across all 39 registered global jurisdictions and custom user-provided datasets.

\end{document}
